\documentclass[11pt,letterpaper]{article}
\usepackage[T1]{fontenc}
\usepackage{lmodern,amsmath,amssymb,amsthm,booktabs,xcolor}
\usepackage[letterpaper,margin=0.91in,top=0.85in,bottom=0.90in,headheight=14pt]{geometry}
\usepackage{fancyhdr,hyperref}
\definecolor{ink}{HTML}{192C40}
\definecolor{teal}{HTML}{096662}
\hypersetup{colorlinks=true,urlcolor=teal,linkcolor=teal,pdftitle={Three transfers from the Dirichlet GRH proof},pdfauthor={Research note prepared for Vladimir Reshetnikov}}
\urlstyle{same}
\setlength{\emergencystretch}{3em}
\newtheorem{theorem}{Theorem}[section]
\pagestyle{fancy}\fancyhf{}
\lhead{\small\sffamily\color{ink}METHODS FROM THE DIRICHLET GRH PROOF}
\rhead{\small\sffamily\color{ink}RESEARCH NOTE}
\cfoot{\small\sffamily\thepage}
\setlength{\parindent}{1.2em}
\setlength{\parskip}{2pt}
\raggedbottom
\newcommand{\ghsource}[2]{\href{https://github.com/lsernau/grh-lean/blob/72fead102f663d5c7a407964a13dc0cf9d977207/#1}{#2}}
\begin{document}
\thispagestyle{plain}
\noindent{\color{teal}\sffamily\bfseries MATHEMATICAL RESEARCH NOTE}\hfill
{\sffamily\small 25 November 2026}\par
\vspace{0.55em}
\noindent{\raggedright\color{ink}\sffamily\bfseries\fontsize{25}{29}\selectfont
Three transfers from the proof\\of Dirichlet GRH\par}
\vspace{0.45em}
\noindent{\large Finite-ring progressions, ramified torsion, and sparse zeta data}\par
\vspace{0.65em}
\noindent\textcolor{teal}{\rule{\linewidth}{0.8pt}}
\paragraph{Results.}
For the dual numbers $R=\mathbb F_2[\epsilon]/(\epsilon^2)$ we construct
a progression-free subset of $R[x]$ of upper density $15/16$.
Every progression-free subset with a natural density has density at most
$119/128$. This gives a strict separation for the full three-term problem
in a case of Asada et al.'s published Question~6.5. We determine the leading
constants for a family of finite rings and both exact extrema for the
stronger condition of avoiding square-ratio endpoint pairs.

For a question posed in the incoming reports, we compute the precise cost
and least degree of a polynomial that blocks every torsion point in a
prescribed order interval on a one-dimensional formal group. Ramification
can put points of different orders at the same distance from zero. The
computation accounts for the distances between points as well. On an
explicit elliptic curve a degree-five obstruction of budget $15/32$ at
order four attains both bounds. We also give formulas along its torsion tower.

Finally we show how sparse data on Li coefficients distinguish Dedekind
zeta functions of arbitrary number fields. Bounded mean absolute discrepancy on even one sequence of increasingly long windows forces their zeta functions to agree. Sparse positivity gives a separate
criterion for the Riemann hypothesis of an individual field; proving the
required signs for a new nonabelian example remains an arithmetic problem.

\paragraph{Reading the source.}
We inspected the shortened Alp\"oge--Claude proof in Luke Sernau's
repository at commit \texttt{72fead102f66}. It proves RH and GRH for Dirichlet
characters of all positive moduli. In the critical strip every zero has
real part $1/2$; in the right half of the strip the equivalent
nonvanishing formulation also extends beyond the strip, with the principal
character at $s=1$ excepted. Its central inequality compares the same
covolume through short integral forms, enough independent forms and a
lower bound from long forms. A one-sided bound on a smoothed prime sum
and a Mellin--Landau argument then exclude off-line zeros.
\ghsource{grh/src/Statements.lean}{Statements};
\ghsource{grh/src/Main.lean}{five theorems};
\ghsource{grh/src/New/Reduction/MainInequality.lean}{covolume};
\ghsource{grh/src/New/Reduction/Zeros/Criterion.lean}{zero criterion}.

The finite-ring argument adapts the source's use of prime exponents to
exclude multiplicative patterns. The formal-group argument takes up its
choice of torsion twists with controlled valuation. The analytic argument
uses the smoothed prime remainder, then transfers the question of what
sampling detects to Li coefficients. The new theorems are paper proofs.
The accompanying computations use exact arithmetic for the finite
examples. We relied on the independent kernel checks for the upstream Lean
result. \href{https://math.mit.edu/~drew/shortgrh/}{Sutherland's verification of both editions}.

\medskip
\noindent\begin{tabular*}{\linewidth}{@{}l@{\extracolsep{\fill}}r@{}}\toprule
\textsf{Results and proofs} & \textsf{Page}\\\midrule
Finite rings and the two density extrema & \pageref{sec:finite-rings}\\
The torsion algorithm and its sharpness proof & \pageref{sec:threshold}\\
A ramified elliptic curve and exact scaling & \pageref{sec:ramified-example}\\
Prime sums, Li coefficients and arithmetic equivalence & \pageref{sec:sampling}\\\bottomrule
\end{tabular*}
\clearpage
\section{Progression-free polynomial sets over finite local rings}
\label{sec:finite-rings}

Asada, Fourakis, Manski, McNew, Miller and Moreland asked in Question~6.5 of
\emph{Subsets of $\mathbb F_q[x]$ free of 3-term geometric progressions}
what changes when the coefficient field is replaced by another finite ring.
We retain their nonunit-ratio convention. The analogue of their degree
ball is the set of polynomials whose coefficients above degree $n$ vanish.
All ratios in this section are polynomials.
\href{https://www.nathanmcnew.com/fqxramsey.pdf}{Author manuscript, pp.~10--12};
\href{https://doi.org/10.1016/j.ffa.2016.10.002}{published paper}.

The Lean connection is an intermediate lemma. In the source,
\texttt{modEq\_orderOf\_of\_div\_pow\_eq} turns a multiplicative equality
into a restriction on exponents, and \texttt{exists\_three\_term\_free}
constructs sets that avoid the resulting additive configurations. An
earlier ProveIt report applied this idea to polynomial valuations over a
field. Over a ring with nilpotents the scalar ratio and its square
act differently. Tracking the prime valuations of the \emph{coefficients
in a basis of the ring} detects forbidden endpoints. The resulting
argument is elementary and independent of the statements about zeros
of $L$-functions.
\href{https://github.com/lsernau/grh-lean/blob/72fead102f663d5c7a407964a13dc0cf9d977207/grh/src/New/Reduction/Data/Existence.lean#L80-L162}{Lean source, lines 80--162};
\href{https://github.com/VladimirReshetnikov/ProveIt/blob/d31039a48d963692e5a5176f47b6603fc9e9d9cf/docs/incoming/grh-progression-densities-and-cubic-prime-race.zip}{earlier field case}.

Let $k=\mathbb F_q$, let $V$ be a $k$-vector space of dimension $b-1$,
and give $R=k\oplus V$ the commutative multiplication
$(u,v)(u',v')=(uu',uv'+u'v)$. Put $Q=|R|=q^b$.
With a fixed basis of $V$, write $f\in R[x]$ as
$(f_0,f_1,\ldots,f_{b-1})$ with entries in $k[x]$.
The ball $B_n$ consists of all such tuples of degrees at most $n$ and has
$Q^{n+1}$ elements. A set is admissible if it contains no triple
$f,rf,r^2f$ with $r\in R[x]$ a nonunit. We work with subsets of
$R[x]\setminus\{0\}$. Retaining zero would change no density.
Write $\Delta(R)$ for the supremum of the upper ball densities of
admissible sets. The natural ball density uses a limit in place of a
limit superior. Reduction of $r=(a,c)$ modulo $V[x]$ detects units:
$r$ is a unit exactly when $a\in k^\times$, as
$(a,c)^{-1}=(a^{-1},-a^{-2}c)$ in this case.

\medskip\noindent\textbf{Theorem (characteristic two).}
Suppose $q$ is even. The set $S$ of nonzero tuples whose monic componentwise
gcd is squarefree is admissible. For every $n\ge2$ it satisfies
\[
 |S\cap B_n|=Q^{n+1}-qQ^{n-1}+q-1.
\]
Its natural density is $1-q/Q^2$. The following bounds include a stronger
construction for upper density:
\begin{equation}
 1-\frac{1}{Q^2}\ \le\ \Delta(R)
 \ \le\ 1-\frac{Q-1}{Q^3-1}
 =1-\frac1{Q^2+Q+1}.
 \label{eq:local-ring-densities}
\end{equation}
For fixed $q$ this determines the leading constant:
$Q^2(1-\Delta(R))\to1$ as $b\to\infty$. All the even-characteristic
assertions hold more generally for a finite commutative local $k$-algebra
$R=k\oplus V$ with nilpotent maximal ideal $V$, provided every $v\in V$
satisfies $v^2=0$. Distinct elements of $V$ may have nonzero product.

\smallskip\noindent\emph{Proof.}
In characteristic two, $r=a+c$ has $r^2=a^2$, since Frobenius
acts coefficientwise on $c\in V[x]$ and all its coefficients square to zero.
In the more general algebra the unit criterion follows by reduction
modulo the nilpotent ideal $V[x]$ and a finite geometric-series inverse. If $a=0$ the last term
of the progression is zero. Otherwise choose a monic irreducible
$P\mid a$. Every component of the last term is divisible by $P^2$.
Conversely every nonzero tuple with a common square factor $P^2$ is a
multiple of the square of the scalar nonunit $P$. Thus $S$ consists
exactly of the nonzero elements divisible by no square of a nonunit.

For any nonzero tuple the indicator that its gcd is squarefree is
$\sum_{d^2\mid\gcd(f_i)}\mu(d)$ over monic $d\in k[x]$. In $B_n$
there are $Q^{n+1-2j}-1$ nonzero multiples of a fixed $d^2$ of degree
$2j\le n$. Unique factorization gives
$\sum_{d\ \mathrm{monic}}\mu(d)z^{\deg d}=1-qz$: it is the reciprocal
of $\sum_{j\ge0}q^j z^j$. Thus only $j=0,1$ contribute after
summing, yielding the count and the asserted natural density.

We next compute the best upper density under the stronger condition that the set
contain no pair $f,a^2f$ with $a\in k[x]$ nonconstant. In a single
ball $B_n$, $n\ge1$, the two layers of highest degrees $n-1,n$ avoid
all such pairs. They have $Q^{n+1}-Q^{n-1}$ elements. For the reverse
bound, map each nonzero $f\in B_n$ to
$x^{2\lfloor(n-\deg f)/2\rfloor}f$ in those two layers. Two distinct
preimages of one point differ by multiplication by a positive even
power of $x$. The map is therefore injective on any pair-avoiding set.

To attain the limiting proportion in a single infinite set, let
$h(f)$ be the largest monic polynomial whose square divides every
component of $f$. In each prime valuation its exponent is half the common
valuation rounded down. Thus
\[
 m(f)=\deg f-2\deg h(f)
\]
is invariant under multiplication by any nonzero scalar polynomial square.
Choose $N_0=-1$ and $N_j=2^{j+1}$ for $j\ge1$. Assign each $f\ne0$
to the unique $j$ with $N_{j-1}<m(f)\le N_j$, and include it exactly
when $\deg f\in\{N_j-1,N_j\}$. All members of a forbidden pair
would have the same assigned $j$ and degrees differing by at least two,
which is impossible. In the top two layers of $B_{N_j}$ a missing
element must have $m(f)\le N_{j-1}$, hence
$\deg h(f)\ge L_j=\lceil(N_j-1-N_{j-1})/2\rceil$.
For any monic $d$ of degree $s\le N_j/2$ the fraction of $B_{N_j}$
divisible componentwise by $d^2$ is $Q^{-2s}$. A union bound
shows that the missed fraction is at most
$\sum_{s\ge L_j}(q/Q^2)^s$, which tends to zero. This proves the
lower bound in \eqref{eq:local-ring-densities} and gives exactly
$1-Q^{-2}$ as the largest attainable upper density under pair avoidance.
The construction also applies for $b=1$.

The natural-density extremum for this same pair condition is also exact.
Each nonzero $f$ factors uniquely as $u h^2$, with $h=h(f)$ monic
and $u$ squarefree in the componentwise gcd sense. For a pair-avoiding
set $E$, let $I_u$ be the set of monic $h$ with $u h^2\in E$.
If $1\in I_u$, there are no other elements. Otherwise
\[
  \sum_{h\in I_u}Q^{-2\deg h}
  \le \sum_{s\ge1}(q/Q^2)^s\le1,
\]
since $q/Q^2\le1/2$. In either case the mass is at most one.
Multiplying by $Q^{-\deg u}$ and summing over base tuples up to degree $n$
gives
\[
  \sum_{f\in E\cap B_n}Q^{-\deg f}
  \ \le\ \sum_{u\in S\cap B_n}Q^{-\deg u}.
\]
The number of all tuples of exact degree $t$ is $(Q-1)Q^t$.
Whenever a set has natural ball density $\delta$, its exact-degree
proportions tend to $\delta$ as well, by subtracting consecutive ball
counts. Division of this last inequality by $(n+1)(Q-1)$ and
convergence of arithmetic means now bound the natural density of $E$ by
$1-q/Q^2$. The squarefree set $S$ attains it. In particular the two
extrema for pair avoidance differ by exactly $(q-1)/Q^2$.

For the upper bound for triple avoidance take any fixed scalar monic linear $P$. Every
nonzero tuple has an exact common $P$-valuation $v$, the minimum of
the component valuations with $v_P(0)=+\infty$. Partition each chain
$h,Ph,P^2h,\ldots$ with $v(h)=0$ into triples starting at exponents
$0,3,6,\ldots$. Multiplication by $P$ is injective and increases the
maximum component degree by one. Hence the triples are disjoint, and
every triple lying within $B_n$ forces at least one deletion. Their
number is exactly
\[
 T_n=\sum_{0\le 3j\le n-2}
 (Q^{n-1-3j}-Q^{n-2-3j}).
\]
Indeed their first terms have degree at most $n-2$ and valuation $3j$.
Dividing by $Q^{n+1}$ and summing the limiting geometric series gives
$(1-Q^{-1})Q^{-2}/(1-Q^{-3})$. This argument works in every
characteristic, including $b=1$. Equation~\eqref{eq:local-ring-densities}
now squeezes $Q^2(1-\Delta(R))$ between $Q^2/(Q^2+Q+1)$ and $1$.\hfill$\square$

\medskip\noindent\textbf{Theorem (a strict gap for triple avoidance).}
Let $D(R)$ be the supremum of natural densities of admissible sets
whose ball densities converge. In the characteristic-two family,
\[
 1-\frac q{Q^2}\le D(R)
 \le 1-\frac q{Q^2}(1-Q^{-1})^q.
\]
Thus as $b\to\infty$ with $q$ fixed, $Q^2(1-D(R))\to q$.
For every $b\ge2$ in this family, $D(R)<\Delta(R)$.

\smallskip\noindent\emph{Proof.}
The lower bound comes from the squarefree construction. To bound natural
density from above, let $C$ consist of nonzero tuples whose gcd has no
linear factor. By reduction modulo the product of the $q$ monic linears,
the natural density of $C$ is $(1-Q^{-1})^q$. For each $g\in C$ form
the $q$ triples $g,Pg,P^2g$, one for each monic linear $P$. Two different
centers give disjoint collections: the common linear valuations of a
member determine the multiplying polynomial $1$, $P$ or $P^2$ uniquely.
Within one collection only the center $g$ is shared. Give a tuple $f$
of degree $d$ the weight $Q^{-d}$. Omitting $g$ costs $Q^{-\deg g}$.
If $g$ is present, at least one of $Pg,P^2g$ is missing for each $P$,
which costs at least $qQ^{-2-\deg g}$. In either case the omitted
weight is at least $qQ^{-2-\deg g}$. Sum over centers in $B_{n-2}$,
so every tuple counted is in $B_n$. Divide by $(n+1)(Q-1)$ and pass to
the limit using the exact-degree proportions of both $C$ and the
omitted set as above. This gives the upper bound. For $b\ge2$ one has
$q(1-Q^{-1})^q>q(1-q/Q)\ge q-1\ge1$, by strict Bernoulli inequality.
Consequently this natural-density upper bound is strictly below the
upper-density construction $1-Q^{-2}$.\hfill$\square$

For $R=\mathbb F_2[\epsilon]/(\epsilon^2)$, the construction has
natural density $7/8$. Every admissible set with a natural density has
density at most $119/128$, while the block construction has upper density
$15/16$. Every admissible subset of $\mathbb F_2[x]$ has
upper density at most $6/7$. More generally any even $q$ and $b\ge2$
give a lower bound at least $1-q^{-3}$, which exceeds the residue-field
upper bound $1-1/(q^2+q+1)$ because $q^3>q^2+q+1$.
The same comparison for an arbitrary $q$ holds in square-zero
extensions of dimension $b\ge3$, by the following calculation.

\medskip\noindent\textbf{Theorem (odd characteristic).}
If $q$ is odd, the nonzero elements divisible by no square of a nonunit
form an admissible set of natural density
\begin{equation}
 \prod_{P\ \mathrm{monic\ irreducible}}
 \bigl(1-z_P^{b+1}+z_P^{b+2}-z_P^{2b+1}\bigr)
 \ \ge\ 1-q^{-b},
 \quad z_P=q^{-\deg P}.
 \label{eq:local-ring-odd}
\end{equation}

\smallskip\noindent\emph{Proof.}
Every forbidden triple ends at a multiple of a nonunit square. Discard
the tuples with $f_0=0$ temporarily, a set of proportion $q^{-(n+1)}$.
We claim the nonzero scalar part $f_0$ of an endpoint forces at least one
prime $P$ to satisfy one of the following two mutually exclusive conditions:
\[
\begin{array}{ll}
 v_P(f_0)=2, &P\mid f_i\quad(1\le i<b),\\
 v_P(f_0)\ge3, &P^2\mid f_i\quad(1\le i<b).
\end{array}
\]
Conversely either condition suffices for square divisibility. To see
necessity write $f=(a,c)^2(h_0,h)=(a^2h_0,a^2h+2ach_0)$
and take $P\mid a$. All nilpotent coordinates are multiples of $P$.
If $v_P(a^2 h_0)\ge3$, either $P^2\mid a$ or $P\mid h_0$,
and all are multiples of $P^2$. To prove sufficiency in the second
case divide by the pure square $P^2$. In the first case set
$h_0=f_0/P^2$. Since both $2$ and $h_0$ are invertible modulo $P$,
choose $c_i$ solving $2c_i h_0\equiv f_i/P\pmod P$. Then
$h_i=(f_i-2Pc_i h_0)/P^2$ is a polynomial for each $i$ and
$f=(P,c)^2(h_0,h)$.

For fixed $P$, uniform residues of independent polynomial coordinates
modulo $P^3$ give total local obstruction probability
$(z_P^2-z_P^3)z_P^{b-1}+z_P^3 z_P^{2(b-1)}$.
The Chinese remainder theorem gives independence for any finite set of
primes once $n$ is sufficiently large. For nonzero $f_0$ an obstruction
with $\deg P=d$ requires $2d\le n$, and has probability at most
$q^{-(b+1)d}$ because it needs $P^2\mid f_0$ and $P\mid f_i$ for all
$i>0$. Thus the omitted large primes have total probability at most
$\sum_{d>D} q^{-bd}$, tending to zero uniformly in $n$.
This gives the product. Since $2b+1\ge b+2$ for $b\ge1$, every
factor is at least $1-z_P^{b+1}$. Their product is $1-q^{-b}$ by
unique factorization.\hfill$\square$

These bounds answer Question~6.5 for an infinite family of finite local
rings. They give both density extrema for endpoint-pair avoidance in characteristic two, and the even
case determines the leading asymptotic of the optimal deficit. The
exact triple-avoidance optimum even for the dual numbers over
$\mathbb F_2$, and comparable results for general finite coefficient
rings remain undetermined here.

\section{The exact cost of blocking a torsion search}
\label{sec:threshold}
In the Lean proof, \texttt{exists\_rootOfUnity\_eval\_le\_gaussVal\_add}
compares a polynomial at a primitive $p$-power root minus one with its Gauss
valuation at that root's radius, with an additive tolerance. It selects a
shell using a harmonic-sum bound and controls contributions from polynomial
roots at that radius. The later simultaneous selection theorem controls
individual polynomials near their respective witness values; the Selmer
argument uses the construction to choose twists. We take up the fixed-radius
variant in the incoming reports.
\ghsource{grh/src/New/Congruences/Selmer/RootsOfUnity.lean\#L292-L305}{Single polynomial};
\ghsource{grh/src/New/Congruences/Selmer/RootsOfUnity.lean\#L474-L486}{simultaneous selection};
\ghsource{grh/src/New/Congruences/Selmer/GeneralTheorem.lean\#L318-L344}{Selmer application}.

One prior report asks for the threshold on a prescribed interval of torsion
orders over a ramified base, including intervals which cannot be
specified by a valuation cutoff. It also asks for the sharp joint dependence
on degree, tolerance and order interval. We answer both questions for
a single polynomial with unrestricted integral coefficients. The input can
in fact be any finite set of points in the open unit disc.
\href{https://github.com/VladimirReshetnikov/ProveIt/blob/d31039a48d963692e5a5176f47b6603fc9e9d9cf/docs/incoming/grh_method_transfers_sharp_torsion.zip}{Prior report, section on formal height and final questions}.

Normalize $v(p)=1$ on $\mathbb C_p$, put $v(0)=+\infty$, and fix a
nonempty finite $S\subset\{z:v(z)>0\}$, a real $\rho>0$ and a tolerance
$\delta\ge0$. The coefficient ring is $\mathcal O_{\mathbb C_p}$.
For a nonzero integral polynomial define its Gauss valuation by
\[
 A(g)=G_\rho(g)=\min_i(v(a_i)+i\rho),\qquad g(T)=\sum_i a_iT^i.
\]
A polynomial \emph{blocks} $S$ when $v(g(z))>A(g)+\delta$ for every $z\in S$.
A point attaining equality is a successful test point.
We compute the smallest possible $A(g)$ of a blocker and its least possible degree.

Partition $S$ by the positive weights $w(z)=\min(\rho,v(z))$ and process
the distinct weights $w$ in decreasing order. Denote the part at $w$ by
$S_w$. On each part take a greedy sequence $z_0,\ldots,z_{l-1}$, $l=|S_w|$:
start at any point, and at step $m<l$ minimize the following quantity over
the unused points $x$:
\[
 E_w(m)=\min_x\sum_{i<m}\bigl(v(x-z_i)-w\bigr),\qquad E_w(l)=+\infty.
\]
In particular $E_w(0)=0$. Given the previously computed integers $m_u$, put
\begin{equation}\label{eq:rec}
 t_w=\delta+\sum_{u>w}m_u(u-w),\qquad
 m_w=\min\{m:E_w(m)>t_w\}.
\end{equation}
The computation uses only the distances in each part and its weight.

\begin{theorem}\label{th:exact}
The least blocking degree $d_*$ and least blocking budget $B_*$ are
\[
 d_*=\sum_w m_w,\qquad B_*=\sum_w w m_w.
\]
One monic polynomial attains both: take the first $m_w$ greedy points of each
part as its roots. Consequently a polynomial with $\deg g\le d$ and
$G_\rho(g)\le B$ is guaranteed a successful test point precisely when
$d<d_*$ or $B<B_*$. With no degree restriction the budget threshold alone
also applies to integral formal power series.
\end{theorem}

\begin{proof}
First consider one part $U=S_w$. Write $P_m=\prod_{i<m}(T-z_i)$ and
$b_m=mw+E_w(m)$ for $0\le m<|U|$. Greedy choice says
$\min_U v(P_m)=v(P_m(z_m))=b_m$. Among all monic polynomials of degree
$m$, this is the largest possible minimum. Indeed write a competing
monic $Q=\sum_{i=0}^m c_i P_i$, with $c_m=1$. If $v(Q(x))>b_m$ for all
$x\in U$, evaluating successively at $z_0,\ldots,z_m$ forces
$v(c_i)+b_i>b_m$ for every $i$. All earlier terms exceed the threshold
everywhere on $U$ and later terms vanish at $z_i$. At $i=m$ we have a
contradiction. Thus tie choices give the same value $E_w(m)$.
At $m=|U|$ the root polynomial vanishes on the whole part.
This is the finite-set argument for Bhargava's $P$-orderings.
\href{https://doi.org/10.1006/jnth.1998.2220}{Bhargava on fixed divisors};
\href{https://arxiv.org/abs/1611.02142}{Amini on tree factorials}.

To see why each part is charged its stated weight, factor a prospective
blocker $g=b\prod_\alpha(T-\alpha)$ in $\mathbb C_p$. Multiplicativity
of the Gauss valuation gives
\[
 G_\rho(g)-G_0(g)=\sum_\alpha
 [\min(\rho,v(\alpha))-\min(0,v(\alpha))],\qquad G_0(g)\ge0.
\]
The contribution of this root to $v(g(z))-G_\rho(g)$ is
$v(z-\alpha)-\min(\rho,v(\alpha))$. Roots with $v(\alpha)\le0$
contribute zero. For any remaining root put $u=\min(\rho,v(\alpha))$.
If there is no part at $u$, the root contributes at most zero at every
test point and can be discarded. If the part exists, its contribution
at every other part of weight $w$ is $\min(u,w)-u$. Within $S_u$ it
comes from the monic product of all roots at weight $u$, followed by
subtraction of $u$ for each root. This reasoning includes the part at
$u=\rho$, which can contain several actual radii or the origin.

Let $M_u$ be the count of retained roots of weight $u$. At each point
of $S_w$ blocking requires the monic product of its $M_w$ roots to have
excess over $M_ww$ strictly greater than
$\delta+\sum_{u>w}M_u(u-w)$. Discarded terms could only have lowered
the excess. By the one-part extremum and induction from the largest
weight down, $M_w\ge m_w$ at every part (the claim is automatic for
$M_w\ge |S_w|$). All retained roots cost their weight in
$G_\rho(g)-G_0(g)$, proving both lower bounds. For the stated monic
polynomial the cross-part effects are exactly the charges in
\eqref{eq:rec}, the one-part minima are attained, and $G_0=0$.
This proves sharpness including both boundary equalities.

For an integral formal series with $G_\rho=A<\infty$, truncate at degree
$D$ with $(D+1)\min_{z\in S}w(z)>A+\delta$. The discarded tail
has valuation above $A+\delta$ at each point and leaves the Gauss
valuation unchanged. The attaining polynomial supplies the reverse direction.
\end{proof}

Updating one running sum per unused point evaluates the greedy sequences
in $O(\sum_w |S_w|^2)$ valuation additions/comparisons given the distances.
The polynomial may require a finite extension of the base field when the
test points are algebraic. Prescribing the field of its coefficients adds
Galois constraints and can change the answer.
For cyclotomic shells of one valuation, earlier reports
\href{https://github.com/VladimirReshetnikov/ProveIt/blob/d31039a48d963692e5a5176f47b6603fc9e9d9cf/docs/incoming/grh_sharp_torsion_and_symplectic_blocks.zip}{already computed a shell recursion}
in terms of the regular cyclotomic tree. Theorem~\ref{th:exact} uses
the actual metric on each selected part, with arbitrary clipping at $\rho$.

\subsection{Obtaining the metric from torsion counts}
Let $F$ be a one-dimensional commutative formal group of finite
reduction height over the integers of a finite extension of $\mathbb Q_p$.
For an interval $h_0\le h\le H$, take $G=F[p^H]$ in the open unit disc and
$D=F[p^{h_0-1}]$. The target is exactly $S=G\setminus D$.
The group law gives $v(x-y)=v(x\mathbin{\ominus_F}y)$: formal subtraction
is divisible by $X-Y$ with unit quotient. Hence $V_r=\{x\in G:v(x)\ge r\}$
is a subgroup and its cosets give the balls at level $r$.

For each distinct torsion valuation, the counts $|V_r|$ and
$|D\cap V_r|$ come from the slopes, with multiplicities, of the
Weierstrass polynomials of $[p^H]$ and $[p^{h_0-1}]$, including the origin
in each subgroup. At the last level use $V_\infty=\{0\}$.
At adjacent levels $r<t$, a coset of $V_r$ has $|V_r|/|V_t|$ children.
If it meets $D$, exactly $|D\cap V_r|/|D\cap V_t|$ children meet $D$,
by translation from a point of intersection. If it misses $D$, all its
children miss $D$. Also distinguish the unique path containing the origin.
This constructs the finite tree up to permutations of equivalent children.
Delete the $D$ leaves. The last common ball of two surviving leaves gives
the valuation of their difference. The analogous comparison with the
marked origin gives each weight. Theorem~\ref{th:exact} now solves the
prescribed order interval using integer populations and rational slopes;
there is no need to solve individually for the torsion coordinates.
For a fixed lower endpoint the first guaranteed upper endpoint, whenever
it exists, is the first $H$ whose computed threshold passes either strict test.
For computation, the input consists of the indicated Newton polygons or
a formal group specified effectively enough to compute them.
\href{https://arxiv.org/abs/math/0604252}{Fargues on ramification filtrations};
\href{https://arxiv.org/abs/0910.3323}{Rabinoff on canonical subgroups}.

For each fixed $B$, increasing $H$ eventually guarantees success even at tolerance zero.
Choose $0<t\le\rho$ smaller than all nonzero valuations in $D$,
with this last condition vacuous for $D=\{0\}$. All torsion points
of valuation $>t$ form a finite subgroup $N$. Indeed,
$[p](X)=pX+O(X^2)$ with integral coefficients raises any valuation
$v>t$ by at least $\min(1,t)$ until the point is killed. A nonzero
point of valuation $>1$ cannot be torsion since from then on the
linear term always uniquely dominates. Thus $N$ is killed by a
bounded power of $p$, and $D\subseteq N\subseteq G$ for large $H$.
Take one representative $x$ of each nonzero coset in $G/N$. With
$c=\sum_{0\ne z\in N}v(z)$ the sum of their valuations is
$C=(H-c)/|N|$: the sum over all nonzero torsion points in $G$ is
$H$, by the derivative $p^H$ at zero and Weierstrass preparation.
The product polynomial $P$ of the representative linear factors
has $G_t(P)=C$ and $v(P'(x))=C-v(x)$, by translation in $G/N$.
All Lagrange polynomials on these nodes have $G_t=0$. If $G_t(g)<C$,
monic division by $P$ leaves an integral remainder with Gauss value
$G_t(g)$, and interpolation gives some node with
$v(g(x))\le G_t(g)\le G_\rho(g)$. Taking $C>B$ proves termination and
explains why the preceding finite algorithm determines the optimal
endpoint as well as testing each prescribed window. This argument uses
the valuation-cutoff interpolation from the prior formal-height report.

\section{An elliptic curve with overlapping torsion valuations}
\label{sec:ramified-example}
Let $K=\mathbb Q_2(\pi)$, with $\pi^4=2$, and let $v(2)=1$. Consider the elliptic curve
\[
                 E:\qquad y^2+\pi xy+y=x^3.
\]
Its discriminant $\pi^3-27$ is a unit. On the special fibre the negation map sends $(x,y)$ to $(x,y+1)$, so the only geometric point killed by $2$ is the identity. The reduction is supersingular. All geometric $2$-power torsion specializes to the identity and is represented in the formal group.

\paragraph{The duplication series.}
Put $t=-x/y$ and $u=-1/y$. The equation becomes $u(1-\pi t-u)=t^3$, so $u=t^3+O(t^4)$ in $\mathcal O_K[[t]]$. Differentiation and substitution into $dx/(2y+\pi x+1)$ give
\[
  \omega=\frac{dt}{1-\pi t-2u}
      =\bigl(1+\pi t+\pi^2t^2+(\pi^3+2)t^3+O(t^4)\bigr)\,dt.
\]
In characteristic zero let $L(t)=\int_0^t\omega$ be the formal logarithm. Comparing coefficients through degree four in $L([2](t))=2L(t)$ gives
\[
                [2](t)=2t-\pi t^2-7t^4+O(t^5).
\]
The coefficient of $t^3$ vanishes; all coefficients of $[2]$ are integral since $E$ has an integral smooth formal group law. These coordinate and invariant-differential constructions are described by Silverman, \emph{The Arithmetic of Elliptic Curves}, second edition, Chapter IV, \S\S1--5, \href{https://doi.org/10.1007/978-0-387-09494-6}{DOI}. For the general theory of subgroups of torsion points closest to zero see Rabinoff, \emph{Higher-level canonical subgroups for $p$-divisible groups}, \href{https://arxiv.org/abs/0910.3323}{arXiv:0910.3323}.

\paragraph{Valuations at every level.}
The negative-slope part of the Newton polygon of $[2](T)$ has vertices $(1,1),(2,1/4),(4,0)$. The three nonzero points killed by $2$ have valuations $3/4$ once and $1/8$ twice. If $c$ denotes the coordinate of the point of valuation $3/4$, its four preimages are roots of $[2](T)-c$, whose negative-slope vertices are $(0,3/4),(2,1/4),(4,0)$. Two preimages have valuation $1/4$ and two have valuation $1/8$. For a nonzero target of valuation $0<y\le 1/2$, the relevant part of the polygon is a single segment from $(0,y)$ to $(4,0)$, possibly containing the point at abscissa $2$. All four preimages have valuation $y/4$. Terms of degree beyond four are integral and do not affect these slopes. Equivalently one can apply Weierstrass preparation after passing to a finite extension containing the parent. The prepared polynomial has degree four, and the remaining unit has constant absolute value one on the open unit disk. The group is in characteristic zero, so preimages are simple.

It follows that for every $n\ge2$ the valuations of the points of exact order $2^n$, and their multiplicities, are given by
\[
\begin{array}{c|ccc}
v(t) & 4^{2-n}/4 & 4^{2-n}/8 & 4^{2-n}/32 \\
\hline
\text{multiplicity} & 2\cdot4^{n-2}&2\cdot4^{n-2}&8\cdot4^{n-2}.
\end{array}
\]
The level $n=2$ follows by lifting the three nonzero points killed by $2$. Induction uses the single-segment calculation, since every valuation on these levels is at most $1/4$. Exact orders double under taking preimages of a point of positive $2$-power order. Levels one and two share the valuation $1/8$. For every $n\ge2$ the last valuation at level $n$ equals the middle valuation at level $n+1$.

\paragraph{The tree for exact order four.}
Let $H=E[4]$, $B=E[2]$, $S=H\setminus B$, and $H_r=\{P\in H:v(t(P))\ge r\}$, with $v(t(O))=+\infty$. Formal subtraction is $(X-Y)$ times an integral power series with constant coefficient one, since it vanishes at $X=Y$ and its linear part is $X-Y$. Thus the ordinary coordinate differences have the same valuation as the formal differences, and the balls in $H$ are cosets of $H_r$. At the successive thresholds,
\[
\begin{array}{c|rrrrr}
r & 1/32 & 1/8 & 1/4 & 3/4 & +\infty\\
\hline
|H_r|       &16&8&4&2&1\\
|B\cap H_r| &4&4&2&2&1 .
\end{array}
\]
Inside the root the two cosets of $H_{1/8}$ contain respectively four and eight members of $S$. In the coset containing $O$, each of the two cosets of $H_{1/4}$ contains two members of $S$, which form one coset of $H_{3/4}$. In the other coset of $H_{1/8}$ all eight points survive. There its two cosets of $H_{1/4}$ each split into two pairs modulo $H_{3/4}$. Each surviving pair splits into two singleton leaves. Consequently the four-point and eight-point branches have different internal structures, even though they are children of the same root.

\subsection{Five roots suffice at every exact order from four on}
Fix $n\ge2$, put $s=4^{2-n}$, and test all points of exact order $2^n$
at tolerance zero and $\rho\ge s/4$. Then the exact answer is
\[
               \boxed{\ d_*=5,\qquad B_*=15s/32.\ }
\]
In particular at order four a budget strictly below $15/32$ ensures a
successful specialization, whereas a quintic of budget exactly $15/32$
blocks every point.

Choose a compatible torsion basis $A_n,B_n$: at level one take the point
of valuation $3/4$ and either point of valuation $1/8$. Lift the first
to valuation $1/4$ at level two and continue by arbitrary preimages.
The basis property follows by reducing a relation repeatedly modulo two,
or by mapping to the two independent basis points of $E[2]$.
Odd multiples preserve valuation by the unit linear coefficient of
odd multiplication in the formal group. On these two cyclic directions
the valuations at exact order $2^j$ are respectively
\[
 a_1=3/4,\quad a_j=4^{1-j}\ (j\ge2),\qquad
 b_j=1/(8\cdot4^{j-1})\ (j\ge1).
\]
The two lists are disjoint. Hence the valuation of $a A_n+b B_n$
is the smaller of the two component valuations, with $+\infty$ for
a zero component. This also computes every pairwise distance by
subtraction of the two coefficient vectors modulo $2^n$.

The inner two parts consist, respectively, of $a$ odd and
$b\equiv0,2\pmod4$. Within either part, differences have $a$ even
and $b$ a multiple of four, hence valuation at least $s/2$
(or $+\infty$). One mark therefore exceeds threshold zero on the
part of weight $s/4$ and threshold $s/8$ on the part of weight $s/8$.
The outer part is $b$ odd, of weight $s/32$. It splits by $b\bmod4$ into
two branches at mutual distance valuation $s/8$. Each splits by
$a\bmod2$, with cross valuation $s/4$. The greedy sequence has
$E(1)=3s/32$, then $E(2)=10s/32$, attained by placing one mark in
each $b$-branch and evaluating at the opposite $a$ parity in a marked
branch. A third mark adds at least $3s/32$ everywhere it matters.
The required excess after the first two parts is precisely $10s/32$.
Its strict inequality therefore costs three outer marks. The total is
$s/4+s/8+3s/32$ as asserted.

\begin{center}
\begin{tabular}{@{}crr@{}}\toprule
Tolerance at exact order four, $\rho=1$ & Minimum degree & Minimum Gauss valuation\\
\midrule
$0$ &5&$15/32$\\
$1/32$ &5&$15/32$\\
$1/8$ &6&$1/2$\\
$1/2$ &9&$29/32$\\
$1$ &12&$1$\\
\bottomrule\end{tabular}
\end{center}
The exact-arithmetic checker evaluates \eqref{eq:rec} in the basis above
and independently enumerates all multisets of marks up to the single-part
size on the order-four test cases. It also verifies the attained
inequality at every test point.

\subsection{Exact scaling in the deep torsion tower}
Suppose the coefficient of $T^Q$ in $[p](T)=\sum_{i\ge1} a_iT^i$ is a unit, where $Q=p^h$, and all earlier coefficients are in the maximal ideal. Put
\[
 c=\min_{1\le i<Q}\frac{v(a_i)}{Q-i}, \qquad D=Qc.
\]
For $z$ in the maximal ideal one has the exact identity
\[
         \min\{v([p](z)),D\}=Q\min\{v(z),c\}.
\]
Indeed, for $v(z)<c$ the term $a_Q z^Q$ has strictly smaller valuation than every other term. For $v(z)\ge c$ every term has valuation at least $D$. This holds also at zero. Formal subtraction therefore gives
\[
 \min\{v(t(x)-t(y)),c\}
  =\frac1Q\min\{v(t([p]x)-t([p]y)),D\}.
\]
Here and below $t(x)$ denotes the coordinate of the formal group point $x$.

Let $S$ be a finite set of nonzero formal torsion points with all $v(t(s))<D$, and let $S'=[p]^{-1}S$. Write $B$ for the optimal obstruction cost on $S$ at tolerance $\delta$ with benchmark $\rho$. If $B+\delta<D$, then the optimal cost on $S'$ at tolerance $\delta/Q$ and benchmark $\rho/Q$ is exactly $B/Q$. The least degree among cost-minimizing obstructions remains the same. When the benchmark exceeds every valuation of $S$, both benchmarks may instead be replaced by one sufficiently large fixed benchmark.

To prove this, use the reduction to multisets of target points. A multiset of marks of cost $b$ defeats a target $s$ exactly when its sum of pair valuations against $s$ exceeds $b+\delta$. For $b\le B$ each valuation in this sum may be capped at $D$ without changing the answer: any summand reaching $D$ already suffices, since all other summands are nonnegative. Every lift has the weight of its parent divided by $Q$, by the first identity above. On $S'$ use the cap $c=D/Q$. The second identity identifies every capped comparison on $S'$ with the corresponding comparison on $S$, divided by $Q$. Sending marks to their parents preserves multiplicities; conversely choose any lift of each occurrence of a mark. Testing the target at all lifts repeats its parent test. Thus both the optimum and the degree comparison follow. Iteration divides cost and tolerance by $Q$ at each step.

For the elliptic curve above, $c=1/8$ and $D=1/2$. At tolerance zero, the exact-order-four cost $15/32$ lies below $D$, so at every exact order $2^n$, $n\ge2$, the cost is $15/(32\cdot4^{n-2})$ with five marks. This argument also applies to overlapping levels. For the union of exact orders $8$ and $16$ the following table gives the shell excesses immediately before and at the first successful mark count $m$. In coordinates $(a,b)\in(\mathbb Z/16)^2$, choose compatible lifts of the closest nonzero $2$-torsion point and of a second independent $2$-torsion point, always taking a lift of maximal valuation. They form a basis as their images under $[8]$ are independent. The valuations of their multiples of order $2^j$ are respectively $(3/4,1/4,1/16,1/64)_j$ and $(1/8,1/32,1/128,1/512)_j$. Multiplication by an odd integer preserves valuation. The lists are disjoint, so the valuation of a sum of the two coordinates is their minimum. In descending order the five shells have congruence descriptions
\[
 (a\equiv2\bmod4,\ b\equiv0\bmod8),\quad
 (a\equiv2\bmod4,\ b\equiv4\bmod8),\quad
 (a\text{ odd},\ 4\mid b),\quad b\equiv2\bmod4,\quad b\text{ odd}.
\]
Each is a translate of a subgroup.
\[
\begin{array}{c|r|c|r|cc}
w& \#S_w& \text{threshold}&m&E_w(m-1)&E_w(m)\\\hline
1/16 &8&0&1&0&1/16\\
1/32 &8&1/32&1&0&3/32\\
1/64 &32&1/16&3&1/16&5/64\\
1/128 &64&13/128&6&13/128&1/8\\
1/512 &128&85/512&12&85/512&25/128
\end{array}
\]
To reproduce the excesses by hand, begin each row at $r_0=w$. For each successive attained finite pair valuation $r_j>w$, let $d_j$ be the number of cosets in that row at threshold $r_j$. At each of these thresholds the balls have equal size. Splitting a prefix evenly among descendants gives $E_w(k)=\sum_{j\ge1}(r_j-r_{j-1})\lfloor k/d_j\rfloor$ for $k<\#S_w$. The values $r_1,r_2,\ldots$ for the first two rows are $1/8,1/4,3/4$. For the third, fourth and fifth rows prepend respectively $(1/32,1/16)$, $(1/64,1/32,1/16)$ and $(1/128,1/64,1/32,1/16)$. In all five cases $d_j=2^{j-1}$. Summing $mw$ gives $27/128$, attained with 23 marks. Its preimages give, for all $N\ge4$, cost $27/(128\cdot4^{N-4})$ with 23 marks on the union of exact orders $2^{N-1}$ and $2^N$.

\paragraph{Eventual scaling for every fixed pattern.}
Fix a benchmark $\rho>0$ and any finite nonempty pattern of $k$ distinct positive exact-order exponents. Translate all these exponents upward by the same integer and take tolerance zero. From an explicit level onward, each unit translation divides the optimal budget by $Q$ and preserves the least degree of a cost-minimizing obstruction. In particular, the budgets are eventually an exact geometric sequence. One sufficient starting condition, writing $j_{\min}$ for the smallest exponent in the translated pattern, is
\[
 D_* = \min(D,\rho),\qquad
 L=1+\left\lfloor\frac{1}{\min(1,D_*)}\right\rfloor,\qquad
 j_{\min}\ \ge\ L+k+2+r,\quad Q^r>k/D_*,\ r\ge0.
\]

To see this, write $H_j=F[p^j]$ and $S_j=H_j\setminus H_{j-1}$. The prepared factorization of $[p^j](T)$ is a unit times $\prod_{y\in H_j}(T-t(y))$. Its coefficient of $T$ has valuation $j$, since $[p^j]'(0)=p^j$. It follows that $\sum_{y\in H_j\setminus\{0\}}v(t(y))=j$, and subtracting successive levels gives $\sum_{y\in S_j}v(t(y))=1$. For any $x\in S_j$ the product of differences against all $y\in H_j$ other than $x$ again has valuation $j$, by translation invariance of the coordinate-distance valuation. Equivalently $[p^j]'(t(x))=p^j$ times a unit, by the invariant differential. Evaluating the prepared factorization of $[p^{j-1}]$ at $x$ then gives the row identity
\[
 \sum_{y\in S_j,\ y\ne x}v(t(x)-t(y))
       =j-v\bigl(t([p^{j-1}]x)\bigr)\ \ge j-1 .
\]
The inequality uses that each nonzero point of $H_1$ has valuation at most its shell's total, namely one.

Every torsion point of valuation at least $D_*$ belongs to $H_L$. Indeed, for a point of valuation $v_0\ge D_*$, a nonzero multiple by $p$ has valuation at least $v_0+\min(1,D_*)$, by the integral multiplication series and its linear term $pT$. Iterating to the last nonzero multiple, which has valuation at most one, gives the stated order bound. Therefore capping all off-diagonal pair valuations in any one row at $D_*$ costs at most $L$ in total: the formal differences are distinct nonzero torsion points, and those whose valuations exceed $D_*$ belong to $H_L$, where the sum of all nonzero valuations is $L$.

Choose a parent pattern whose smallest exponent is at least $L+k+2$. All its point valuations are strictly below $D_*$. Mark every parent point once; the cost is $k$. At any parent target, just its own exact-order shell contributes off-diagonal capped row sum at least $j-1-L>k$. Now lift the pattern by $r$ exponents and choose one arbitrary lift under $[p^r]$ for each marked parent. Iterating the capped identity (which also holds with $D$ replaced by any smaller positive cap) gives, for all points $x,y$,
\[
 \min\{v(t([p^r]x)-t([p^r]y)),D_*\}
 =Q^r\min\{v(t(x)-t(y)),D_*/Q^r\}.
\]
Each chosen mark has weight equal to its parent's weight divided by $Q^r$. Every target in the lifted pattern lies over a parent target, so its capped row sum exceeds $k/Q^r$. This constructs an obstruction of cost $k/Q^r<D_*$. Its optimal cost is also below $D_*$, and the exact scaling result now applies at each further level. All these target and mark valuations stay below $\rho$, as required for the fixed benchmark.

\section{Extremal prime sums on a prescribed sparse sequence}\label{sec:sampling}

The earlier ProveIt note \emph{Torsion and prime sampling} established a
one-sided zero criterion on prescribed power grids and on much thinner
grids for a residual set of bases. It asked how much survives for a
prescribed base. Here one and the same explicit set of integer exponents
works for every real base $b>1$, including $b=2$. It recovers the exact
upper limit under critical-line RH and detects an attained rightmost
zero without any zero hypothesis. The set of exponents is
\[
 M_j=2^{2^{j^2}},\qquad
 \mathcal E=\bigcup_{j\ge1}\{M_j,2M_j,\ldots,jM_j\}.
\]
Among the cutoffs $x=b^e$, $e\in\mathcal E$, at most
$O_b(1+\log\log\log X)$ lie below $X$ as $X\to\infty$.
Indeed $b^{M_j}\le X$ implies $j^2=O_b(1+\log\log\log X)$,
and the first $j$ blocks contain at most $j(j+1)/2$ terms.

For a number field $K$, write $\Lambda_K(n)$ for the coefficients of
$-\zeta_K'/\zeta_K$ and $\gamma_K$ for its Euler--Kronecker constant,
with the convention $\zeta_K'/\zeta_K(s)=-(s-1)^{-1}+\gamma_K+O(s-1)$.
Take the ramp used in the preceding note, modeled on the \ghsource{grh/src/New/Reduction/OneSidedBound.lean}{prime remainder in the Lean proof},
\[
 R_K(x)=\sum_{n\le x}\Lambda_K(n)(1/n-1/x)-\log x+1+\gamma_K.
\]
For each nontrivial zero $\rho=\sigma+i\gamma$ put
$c_\rho=1/[\rho(1-\rho)]$ and let $\Theta$ be the supremum of the
real parts of all these zeros. All sums count multiplicity.

\paragraph{Theorem.}
If every nontrivial zero of $\zeta_K$ has real part $1/2$, then
for every $b>1$,
\[
 \lim_{j\to\infty}\max_{1\le k\le j}
 b^{kM_j/2}R_K(b^{kM_j})
 = \sum_{\rho}\frac{1}{1/4+(\operatorname{Im}\rho)^2}
 =: S_K.
\]
For $K=\mathbb Q$ the repository's RH theorem applies and
$S_{\mathbb Q}=2+\gamma_{\mathbb Q}-\log(4\pi)
=0.0461914\ldots$.
Unconditionally for every number field, the limsup of
$x^{1-\Theta}R_K(x)$ on $x=b^e$, $e\in\mathcal E$, is strictly positive
if and only if some zero has real part $\Theta$. When no zero attains
$\Theta$, that normalized expression tends to zero on the entire ray.

\paragraph{Proof.}
The smoothed explicit formula is
\[
 R_K(x)=\sum_\rho c_\rho x^{\rho-1}
       +O_K(x^{-1}\log(2x)),\qquad x\ge2 .
\]
For the normalization, the uncentered sum is $(1-1/x)\Phi_K(x)$
in Ihara's notation. His explicit formula (1.2.1), the archimedean
term (1.2.4), and the constant identity (1.3.11) give the displayed
formula. Equivalently, multiplying $-\zeta_K'/\zeta_K(s)$ by
$x^{s-1}/[s(s-1)]$ gives the residues $c_\rho x^{\rho-1}$
and $\log x-1-\gamma_K$ at the zeros and the pole at $1$,
respectively. The zero count $O_K(T\log(2T))$ implies
$\sum_\rho|c_\rho|<\infty$.
By dominated convergence, the sum over $\operatorname{Re}\rho<\Theta$
is $o(1)$ after multiplying by $x^{1-\Theta}$. The normalized
error also tends to zero. Thus the only surviving sum is
$\sum_{\operatorname{Re}\rho=\Theta}c_\rho
\exp(i\gamma\log x)$.

Fix any finite list of its $d$ ordinates, set
$\theta_i=\gamma_i\log b/(2\pi)$, and take an integer $Q\ge2$.
The fractional parts of $0,M_j\theta,\ldots,Q^dM_j\theta$ occupy
$Q^d$ half-open boxes of side $1/Q$ in $[0,1)^d$. The difference between
two points in one box supplies $1\le k\le Q^d$ with
$\|kM_j\theta_i\|_{\mathbb R/\mathbb Z}\le1/Q$ for all $i$.
For every $j\ge Q^d$ this $k$ belongs to the specified block.
In fact a positive fraction of each large block works. For $N=j+1$
points and $B=Q^d$ boxes, the number of same-box pairs is at least
$(N^2/B-N)/2$ by Cauchy--Schwarz. Each positive difference occurs
in at most $N$ pairs. Thus there are at least $(N/B-1)/2$ distinct
suitable $k$ in the block.
Absolute summability lets us first discard a tail of arbitrarily
small norm and then choose $Q$ to make all remaining phases arbitrarily
close to one. Hence a subsequence of the normalized values tends to
$\sum_{\operatorname{Re}\rho=\Theta}c_\rho$. This sum is real by
conjugation. Every summand has positive real part because
$\operatorname{Re}[\rho(1-\rho)]=\sigma(1-\sigma)+\gamma^2>0$.
This proves the assertion about attainment, including the empty case.
On the critical line every $c_\rho$ is positive real. The triangle
inequality supplies the matching upper bound on every block. The
value of the zero sum is explicit in field invariants. For signature
$(r_1,r_2)$, discriminant $d_K$ and Euler's constant $\gamma$,
Ihara (1.2.2), (1.2.3) and (1.3.11) give
\[
 S_K=2+2\gamma_K+\log|d_K|-r_1(\gamma+\log(4\pi))
                    -2r_2(\gamma+\log(2\pi)).
\]
Taking $K=\mathbb Q$ gives the stated numerical constant. Moreover,
for every fixed $\varepsilon>0$, a positive fraction of every sufficiently
large block comes within $\varepsilon$ of $S_K$ on RH. This follows by
fixing the tail and $Q$ for that tolerance in the same-box count.

\paragraph{Consequence and scope.}
For an arbitrary field $K$, a finite upper bound on
$\sqrt{x}R_K(x)$ on this grid rules out an attained rightmost zero
with real part greater than $1/2$. Indeed, on the subsequence just
constructed, the extra factor $x^{\Theta-1/2}$ tends to infinity.
The case where off-line real parts approach their supremum strictly
from below remains open for this test. The earlier question also specified
a different, preassigned sequence of exponents $\exp(n/\log(n+2))$.
The block recurrence proved here addresses the fixed-base issue on
the displayed grid. Sparsity counts cutoffs and conveys no
computational savings in evaluating a prime sum at such a large cutoff.

\smallskip\noindent
For the explicit formula see Y.~Ihara, \href{https://www.kurims.kyoto-u.ac.jp/~kenkyubu/emeritus/ihara/Publications-and-Recent-Preprints/RecentArticles/pdf-files/1A.pdf}{\emph{On the Euler--Kronecker constants of global fields and primes with small norms}},
Part~1, \S\S1.2--1.3.
The antecedent ProveIt argument and its open question appear in
\href{https://github.com/VladimirReshetnikov/ProveIt/blob/d31039a48d963692e5a5176f47b6603fc9e9d9cf/docs/incoming/grh-torsion-and-prime-sampling.zip}{\texttt{grh-torsion-and-prime-sampling.zip}},
in the section on sparse scales.

\subsection{A sparse Li criterion for every number field}

The preceding prime-sum test loses off-line zeros when their real parts
approach an unattained supremum. The same block argument gives a full
equivalence for a different arithmetic sequence. For a number field $K$
set $w_\rho=1-1/\rho$ and define its Li coefficients by
\[
 \lambda_n(K)=\lim_{T\to\infty}\sum_{|\operatorname{Im}\rho|\le T}
            (1-w_\rho^n),\qquad n\ge1.
\]
The sums are real by conjugation. Equivalently, if $\xi_K$ is the entire
completed Dedekind zeta function with its poles removed, then
$\lambda_n(K)=\frac{1}{(n-1)!}(d/ds)^n[s^{n-1}\log\xi_K(s)]_{s=1}$.

\paragraph{Corollary.}
For every number field $K$, critical-line RH for $\zeta_K$ is equivalent
to $\lambda_n(K)\ge0$ for all sufficiently large $n\in\mathcal E$.
The same equivalence holds if the fraction of negative coefficients
in the $j$th block tends to zero. The index set samples at most
$O(1+\log\log N)$ coefficients of index at most $N$.

\paragraph{Proof.}
On RH, $|w_\rho|=1$ and every real part $1-\operatorname{Re}(w_\rho^n)$
is nonnegative. If RH fails, symmetry about the critical line gives
a zero with real part below $1/2$, hence $|w_\rho|>1$ since
$|w_\rho|^2=1+(1-2\operatorname{Re}\rho)/|\rho|^2$. The numbers
$w_\rho$ tend to $1$ at increasing height. Consequently the largest
modulus $R>1$ is attained at finitely many zeros, say $m$ with
multiplicity, and all other moduli are at most $r$ for some $1<r<R$.
For all these other zeros Taylor's formula on the segment from $1$ to
$w_\rho$ gives, for $n\ge2$,
\[
 \left|1-w_\rho^n-\frac n\rho\right|
 \le \frac{n(n-1)}{2|\rho|^2}r^{n-2}.
\]
The sum of $|\rho|^{-2}$ is finite by the zero count. The symmetric sum
of $1/\rho$ exists since paired conjugates have sum
$2\operatorname{Re}\rho/|\rho|^2$. It follows that
\[
 R^{-n}\lambda_n(K)
 =-\sum_{|w_\rho|=R}(w_\rho/R)^n+o(1).
\]
Apply the same-box argument to the arguments of the $m$ outermost
points. The real parts of their normalized powers can all be made
arbitrarily close to $1$ in every sufficiently large block, and the
real part of their sum never exceeds $m$. Thus
\[
 \lim_{j\to\infty}\min_{1\le k\le j}
 R^{-kM_j}\lambda_{kM_j}(K)=-m.
\]
Taking the phase tolerance fixed and small shows that a positive
fraction of every large block has negative coefficients if RH fails.
Finally, counting blocks with $M_j\le N$ gives the claimed sparsity.

Li's number-field criterion and the dominant exponential terms in
its failure case are classical. The additional conclusion is the
prescribed sparse set of indices, its robustness to a vanishing
fraction of exceptions in each block, and the normalized block limit.
In fact the proof permits any sequence of positive integer dilations
$M_j\to\infty$ with block lengths tending to infinity. For
$K=\mathbb Q(\sqrt[3]{2})$, for example, proving the sampled
inequalities would establish RH for that field, which includes a
nonabelian factor. We have derived the equivalence. An arithmetic proof
of the inequalities for this field remains open.
See E.~Bombieri and J.~C.~Lagarias,
\href{https://websites.umich.edu/lagarias/doc/bombieri.pdf}{\emph{Complements to Li's
criterion for the Riemann hypothesis}}, Theorem~1, and
J.~C.~Lagarias,
\href{https://doi.org/10.5802/aif.2311}{\emph{Li coefficients for automorphic
$L$-functions}}, especially the introduction on the Dedekind case.

\subsection{Arithmetic equivalence from sparse Li coefficients}

Let $d_n=\lambda_n(K)-\lambda_n(L)$ for two number fields $K,L$.
The following inverse result requires no critical-line hypothesis.

\paragraph{Theorem.}
If $K,L$ are arithmetically inequivalent, meaning that
$\zeta_K\ne\zeta_L$, then
\[
 \lim_{H\to\infty}\ \inf_{A\ge1}
   \frac1H\sum_{n=A}^{A+H-1}|d_n|^2=+\infty,
 \qquad A,H\in\mathbb N .
\]
Thus even a bounded mean square on one sequence of windows of increasing
length forces $\zeta_K=\zeta_L$. The windows may
be arbitrarily far apart. For example the windows starting at
$A_j=2^{2^{j^2}}$ and of length $j$ use only $O(1+\log\log N)$ indices
up to $N$. Equality of zeta functions gives $d_n=0$ for every $n$.

\paragraph{Proof.}
Write $m_w$ for the signed difference of the multiplicities of the
transformed zero $w=1-1/\rho$. Taking second differences of the two
symmetric zero sums gives
\[
 a_n:=-(d_{n+2}-2d_{n+1}+d_n)=\sum_w c_w w^n,\qquad
 c_w=m_w(1-w)^2,\qquad \sum_w|c_w|<\infty.
\]
Indeed, $|1-w|=|\rho|^{-1}$ and the zero count gives convergence
of the square sum. All $|w|$ are bounded. Suppose, towards the
contrapositive, that there are windows of length $H\to\infty$
and start $A$ on which the mean square of $d_n$ is at most $C^2$.
Then $|d_n|\le C\sqrt H$ throughout each window.

If an unmatched $w$ has modulus above $1$, let $R$ be the largest
such modulus. It is attained, since $w\to1$ at increasing height.
There are finitely many distinct unmatched points $w_1,\ldots,w_t$
with this modulus, and the other unmatched moduli are bounded by
some $r<R$. Put $z_i=w_i/R$. For $N=A+H-t-2$ and $0\le h<t$,
\[
 R^{-N-h}a_{N+h}
 =\sum_{i=1}^t z_i^h c_{w_i}z_i^N +O((r/R)^{N+h}).
\]
The left sides tend to zero since $N\ge H-t-1$ and
$|a_{N+h}|\le4C\sqrt H$. The fixed Vandermonde matrix
$(z_i^h)_{h,i}$ is invertible. Consequently $c_{w_i}z_i^N\to0$,
contradicting $|c_{w_i}|>0$. Thus all outer points match.
Reflection $\rho\mapsto1-\rho$ induces $w\mapsto1/w$ and preserves
multiplicity, so all inner points match as well.

All remaining terms have $|w|=1$. Fix such a $z$.
The average of $a_n z^{-n}$ over $A\le n\le A+H-3$, normalized
by $H$, tends to $c_z$. Indeed, each other mass has geometric
average bounded in modulus by $\min(1,2/(H|1-w/z|))$, uniformly
in $A$, and the masses are absolutely summable. Shifting the two
indices in the second difference now gives
\[
 \frac1H\sum_{n=A}^{A+H-3}a_nz^{-n}
 =-(z-1)^2\frac1H\sum_{n=A}^{A+H-1}d_nz^{-n}
       +O(C/\sqrt H).
\]
Every omitted endpoint is bounded by $C\sqrt H$. Since $z\ne1$,
the normalized Fourier averages of $d_n$ tend to $-m_z$.
For any finite set of distinct unit points, their exponentials on
the window have normalized Gram matrix tending to the identity, by
the same geometric-sum estimate. The projection inequality
(or Cauchy--Schwarz applied to their linear combinations) therefore
gives $\sum_z |m_z|^2\le C^2$ on every such finite set.
The multiplicities are integers. Only finitely many can differ.

The pole-cleared completions $\xi_K,\xi_L$ are entire of order one.
Hadamard factorization gives $\xi_K/\xi_L=e^{a+bs}P(s)/Q(s)$
for polynomials $P,Q$. Set $Z(s)=\zeta_K(s)/\zeta_L(s)$.
The gamma factors of a number field are products of
$\Gamma_{\mathbb R}(s)$ and $\Gamma_{\mathbb R}(s+1)$ with
$\Gamma_{\mathbb R}(s)=\pi^{-s/2}\Gamma(s/2)$. Their two-step
quotients are rational functions, since
$\Gamma_{\mathbb R}(s+2)/\Gamma_{\mathbb R}(s)=s/(2\pi)$.
The discriminant exponential contributes a constant under the same
shift. Hence $Z(s+2)/Z(s)$ is rational. On the positive real axis,
$\zeta_K(s)=1+O_K(2^{-s})$ by absolute convergence of the
Dirichlet series at $2$, and likewise for $L$. Thus
$Z(s+2)/Z(s)=1+O_{K,L}(2^{-s})$. A rational function with this
decay towards $1$ must be identically $1$. For real $s>1$,
$Z(s)=Z(s+2m)\to1$, proving equality of the two zeta functions.

This identifies a zeta function up to equality and allows nonisomorphic
arithmetically equivalent fields. The theorem uses the classical Li
zero sum and a uniform Fourier average on long integer windows.
It supplies an infinite-data test and gives no effective length at
which a specified pair of fields can be distinguished.

\paragraph{Strengthening to mean absolute error.}
In the rigidity theorem one may replace $|d_n|^2$ by $|d_n|$.
Here is a proof using the idempotent theorem in harmonic analysis.
Suppose the mean absolute error is bounded by $C$ on windows of
length $H\to\infty$. Each $|d_n|$ on such a window is at most $CH$,
so the preceding Vandermonde argument, applied near its end, still
eliminates unmatched outer and inner points. Let $\Gamma$ be the
countable subgroup of the circle generated by the remaining points $w$,
given the discrete topology. Its dual $G$ is compact and metrizable.
Write $g_n\in G$ for the character $w\mapsto w^n$.

Take triangular nonnegative weights $t_n$ of total mass one on the
window with its last two positions omitted. Explicitly, normalize the
sequence $\min(i+1,H-2-i)$ for $0\le i\le H-3$, with zero values
outside. They satisfy $\sup t_n=O(1/H)$ and
$\sup|t_{n+1}-t_n|=O(1/H^2)$. The measures
$\sum_n t_nd_n\delta_{g_n}$ have bounded variation, so converge
weak-star on a subsequence to a finite measure $\nu$.
For each $z\in\Gamma$, Abel summation shows that the $t_n$-weighted
average of $(w/z)^n$ tends to zero for $w\ne z$ (the total variation
of the weights is $O(1/H)$). The weighted average of $a_nz^{-n}$
therefore tends to $c_z$, taking $c_z=0$ for points outside the
spectrum. Shifting indices as above introduces an error at most
$O(H^{-2}\sum_{n\in I}|d_n|)=O(C/H)$. Hence
$\widehat\nu(z)=-m_z$ for $z\ne1$. Subtracting
$\widehat\nu(1)$ times Haar probability sets the trivial coefficient
to zero.

The resulting Fourier transform is integer-valued and bounded.
Choose a polynomial taking value zero at zero and one at every
nonzero integer in its finite range. Applying its convolution
polynomial to the measure gives a finite measure with Fourier
transform $1_D$, where $D=\{w:m_w\ne0\}$. The
Cohen--Helson--Rudin idempotent theorem puts $D$ in the coset ring
of $\Gamma$. See B.~Green and T.~Sanders,
\href{https://doi.org/10.4007/annals.2008.168.1025}{\emph{A quantitative version of the
idempotent theorem in harmonic analysis}}.

An infinite coset-ring subset of a subgroup of the circle has closure
the entire circle in its usual topology. To verify this, express a
Boolean combination of finitely many cosets as a union of atoms.
Intersecting the positive cosets in one atom yields a single coset
of a subgroup $T$ (possibly the whole group). Split it into finitely
many smaller cosets to settle all the excluded cosets whose subgroups
intersect $T$ in finite index, using the intersection of those
finite-index subgroups. Each surviving piece is a coset of a subgroup
$T_0$ with finitely many cosets of infinite-index subgroups removed.
If $T_0$ is infinite it is dense in the circle. In a
translation-invariant finitely additive probability on the abelian
group $T_0$, every infinite-index coset has measure zero by disjoint
translates. Every nonempty open arc intersected with $T_0$ has positive
measure, since finitely many translates cover $T_0$. Each such
surviving infinite piece is therefore dense. But the only possible
accumulation point of $D$ is $1$, since $w_\rho\to1$ as the zero
height grows. Thus $D$ is finite. The same Hadamard and shift argument
now gives $\zeta_K=\zeta_L$.


\clearpage
\section{Further research questions and evidence}
\paragraph{Finite rings.}
For the dual numbers over $\mathbb F_2$, the full progression problem
still leaves an interval for each optimum. Does allowing both endpoints
of a square-ratio pair increase the achievable upper density beyond
$15/16$? In odd characteristic the product in \eqref{eq:local-ring-odd}
gives an explicit construction. Determining the leading constant in
that setting would clarify the role of the mixed term in $(a+c)^2$.
The field conjecture numbered 6.3 by Asada et al. concerns
$\mathbb F_q[x]$ and remains separate from their finite-ring question.

\paragraph{Torsion.}
The extremizing roots in Theorem~\ref{th:exact} lie among the tested
geometric points. If the polynomial must have coefficients in the
elliptic curve's field of definition, which of the exact thresholds can
still be attained? The unrestricted construction allows individual
points to be selected even when their Galois conjugates are omitted.
For simultaneous tests on several polynomials an additional choice is
whether to bound a sum of valuations or each positive increase individually.

\paragraph{Sparse analytic data.}
The positivity criterion supplies a definite target for
$K=\mathbb Q(\sqrt[3]{2})$. The original covolume calculation is tied to
Dirichlet twists and a specific family of theta lifts; this report does
not supply the corresponding arithmetic inequality for $K$.
The inverse theorem asks for infinitely many samples in windows of
growing length. Effective finite bounds would require quantitative
control of zero ordinates, separation and tails as the fields vary.

\paragraph{Provenance and reproduction.}
The included inventory identifies all 35 ZIPs in the incoming directory
at ProveIt commit \texttt{d31039a48d96}. The closest antecedents are the
progression-density paper, the sharp-torsion reports, and the paper on
torsion and prime sampling. Their specialized conclusions are identified
where used. The published fixed-divisor, formal-group, explicit-formula
and Li-coefficient ingredients are cited in the proofs. Priority for the
precise combinations established here is provisional.
The archive README gives the build commands and the inputs and outputs
of the finite checks.
\begin{center}
\begin{tabular}{@{}p{0.32\linewidth}p{0.58\linewidth}@{}}\toprule
\textsf{Check} & \textsf{What was compared}\\\midrule
Finite-ring arithmetic & Direct enumeration of gcds, local square multiples,
collections of disjoint triples and pairs across degree layers.\\[3pt]
One-part torsion minima & Greedy fixed divisors against exhaustive multisets
allowing repeated roots in the small example.\\[3pt]
Torsion windows & Coordinate differences in an elliptic basis against a
tree reconstructed from subgroup counts in 20 parameter choices.\\[3pt]
Two-level obstruction & An explicit list of 23 roots and the five intermediate
thresholds at orders eight and sixteen.\\\bottomrule
\end{tabular}
\end{center}
\end{document}
